\subsection{Hermitian and Normal Matrix Structure}

\begin{theorem} \label{thm:simultaneous-congruence-one-positive-definite-form} \hlt{Simultaneous Congruence with One Positive Definite Form}\\
Let $n\in\mathbb{Z}_{>0}$, $\mathbb{F}\in\{\mathbb{R},\mathbb{C}\}$, and let $A,B\in\operatorname{Mat}(n,n,\mathbb{F})$ satisfy $A=A^*$, $B=B^*$, and $A\succ0$.\\
Then $(\exists\,S\in\operatorname{Mat}(n,n,\mathbb{F})^\times)(\exists\,\lambda_1,\ldots,\lambda_n\in\mathbb{R})$ such that, for $\Lambda\defeq\operatorname{diag}(\lambda_1,\ldots,\lambda_n)$:
\begin{enumerate}[label=(\roman*)]
\setlength{\itemsep}{0pt}
\item $A=SS^*$ and $B=S\Lambda S^*$.
\item $p_{A^{-1}B}(z)=\prod\limits_{j=1}^{n}(z-\lambda_j)$.
\end{enumerate}
\begin{proof}
Choose an invertible $T$ with $A=TT^*$ by Theorem \ref{thm:positive-definite-matrix-equivalences}(iv), and set $C\defeq T^{-1}B(T^*)^{-1}$.\\
The identity $C^*=C$ holds by Propositions \ref{prop:basic-adjoint-identities}(ii),(iii) and \ref{prop:adjoints-inverses-and-polynomials}(i). Thus $C$ is self-adjoint and normal.\\
If $\mathbb{F}=\mathbb{R}$, choose a unitary $U$ and diagonal $\Lambda$ with $C=U\Lambda U^*$ by Theorem \ref{thm:real-spectral-theorem}(i)$\Rightarrow$(iii).\\
If $\mathbb{F}=\mathbb{C}$, make the same choice by Theorem \ref{thm:complex-spectral-theorem}(i)$\Rightarrow$(iii).\\
Its diagonal entries are real automatically over $\mathbb{R}$ and by Theorem \ref{thm:self-adjoint-eigenvalues-real} over $\mathbb{C}$.\\
Write $\Lambda=\operatorname{diag}(\lambda_1,\ldots,\lambda_n)$ with $\lambda_1,\ldots,\lambda_n\in\mathbb{R}$, and set $S\defeq TU$.
\begin{enumerate}[label=(\roman*)]
\setlength{\itemsep}{0pt}
\item The matrix $S$ is invertible, and $SS^*=TUU^*T^*=A$ and $S\Lambda S^*=TCT^*=B$.
\item The identity $A^{-1}B=(S^*)^{-1}\Lambda S^*$ holds by part (i).\\
Hence $p_{A^{-1}B}=p_{\Lambda}$ by Theorem \ref{thm:similarity-preserves-characteristic-polynomial}(i). The identity $p_{\Lambda}(z)=\prod\limits_{j=1}^{n}(z-\lambda_j)$ holds by Proposition \ref{prop:upper-triangular-characteristic-polynomial}(i).
\end{enumerate}
\hlt{Structural Commentary}: The Gram factor turns simultaneous congruence into one self-adjoint diagonalization.\\
\hlt{Key Insight}: Diagonalizing $C$ once produces both $S$ and the generalized eigenvalues.
\end{proof}
\end{theorem}

\begin{corollary} \label{cor:positive-definite-hermitian-product-spectrum} \hlt{Products with One Positive Definite Factor}\\
Let $n\in\mathbb{Z}_{>0}$, $\mathbb{F}\in\{\mathbb{R},\mathbb{C}\}$, and let $A,B\in\operatorname{Mat}(n,n,\mathbb{F})$ satisfy $A\succ0$.
\begin{enumerate}[label=(\roman*)]
\setlength{\itemsep}{0pt}
\item If $B=B^*$, then $AB$ is diagonalizable over $\mathbb{F}$ and $\operatorname{spec}(AB)\subseteq\mathbb{R}$.
\item If $B\succeq0$, then $AB$ is diagonalizable over $\mathbb{F}$ and $\operatorname{spec}(AB)\subseteq[0,\infty)$.
\item If $B\succeq0$ and $B\neq0$, then $(\exists\,\lambda\in\operatorname{spec}(AB))\,\lambda>0$.
\item If $B\succ0$, then $AB$ is diagonalizable over $\mathbb{F}$ and $\operatorname{spec}(AB)\subseteq(0,\infty)$.
\end{enumerate}
\begin{proof}
The matrix $A^{-1}\succ0$ by Theorem \ref{thm:positive-definite-matrix-equivalences}(i)$\Rightarrow$(v). If $B=B^*$, choose $S$ and $\Lambda$ by Theorem \ref{thm:simultaneous-congruence-one-positive-definite-form} for $(A^{-1},B)$.\\
Write $\Lambda=\operatorname{diag}(\lambda_1,\ldots,\lambda_n)$, where $\lambda_1,\ldots,\lambda_n\in\mathbb{R}$, $S$ invertible. Then $A^{-1}=SS^*$ and $B=S\Lambda S^*$, so $AB=(S^*)^{-1}\Lambda S^*$.
\begin{enumerate}[label=(\roman*)]
\setlength{\itemsep}{0pt}
\item The identity $AB=(S^*)^{-1}\Lambda S^*$ diagonalizes $AB$ over $\mathbb{F}$ with real eigenvalues $\lambda_1,\ldots,\lambda_n$.
\item The identity $B=B^*$ holds by positive semidefiniteness, so part (i) applies.\\
Set $\mathbf{x}_j\defeq(S^*)^{-1}\mathbf{e}_j$ for $j\in\{1,\ldots,n\}$. Then $(\forall\,j\in\{1,\ldots,n\})\,\lambda_j=\mathbf{x}_j^*B\mathbf{x}_j\geq0$.
\item By parts (i) and (ii), $\operatorname{spec}(AB)=\{\lambda_1,\ldots,\lambda_n\}\subseteq[0,\infty)$. Since $B=S\Lambda S^*\neq0$, one has $\Lambda\neq0$, so $(\exists\,\lambda\in\operatorname{spec}(AB))\,\lambda>0$.
\item The identity $B=B^*$ holds by positive definiteness, so part (i) applies. Set $\mathbf{x}_j\defeq(S^*)^{-1}\mathbf{e}_j$ for $j\in\{1,\ldots,n\}$.\\
Since $S$ is invertible, $(\forall\,j\in\{1,\ldots,n\})[\mathbf{x}_j\neq\mathbf{0}\land\lambda_j=\mathbf{x}_j^*B\mathbf{x}_j>0]$.
\end{enumerate}
\hlt{Structural Commentary}: The simultaneous form of $(A^{-1},B)$ makes $AB$ similar to the same real diagonal matrix.\\
\hlt{Key Insight}: Its eigenvalue signs are the quadratic-form signs of $B$ on transformed coordinate vectors.
\end{proof}
\end{corollary}

\begin{definition} \label{def:phase-inertia-matrix} \hlt{Phase Inertia Matrix}\\
Let $n\in\mathbb{Z}_{>0}$ and $D=[d_{jk}]\in\operatorname{Mat}(n,n,\mathbb{C})$. The matrix $D$ is a \hlt{phase inertia matrix} iff $D$ is diagonal and $(\forall\,j\in\{1,\ldots,n\})\,\abs{d_{jj}}\in\{0,1\}$.
\end{definition}

\begin{proposition} \label{prop:normal-matrix-phase-inertia-form} \hlt{Phase Inertia Form for Normal Matrices}\\
Let $n\in\mathbb{Z}_{>0}$ and let $A\in\operatorname{Mat}(n,n,\mathbb{C})$ be normal. Let $\lambda_1,\ldots,\lambda_n$ list its eigenvalues with algebraic multiplicity.\\
For $j\in\{1,\ldots,n\}$, set $d_j\defeq 0$ if $\lambda_j=0$ and $d_j\defeq\frac{\lambda_j}{\abs{\lambda_j}}$ otherwise. Put $D\defeq\operatorname{diag}(d_1,\ldots,d_n)$.
\begin{enumerate}[label=(\roman*)]
\setlength{\itemsep}{0pt}
\item $D$ is a phase inertia matrix.
\item $(\exists\,S\in\operatorname{Mat}(n,n,\mathbb{C})^\times)\,A=SDS^*$.
\end{enumerate}
\begin{proof}
\begin{enumerate}[label=(\roman*)]
\setlength{\itemsep}{0pt}
\item By the definition of $d_j$, $(\forall\,j\in\{1,\ldots,n\})\,\abs{d_j}\in\{0,1\}$. Hence $D$ is a phase inertia matrix.
\item Set $L\defeq\operatorname{diag}(\lambda_1,\ldots,\lambda_n)$. After permuting its eigenbasis, choose a unitary $U$ with $A=ULU^*$ by Theorem \ref{thm:complex-spectral-theorem}(i)$\Rightarrow$(iii).\\
For $j\in\{1,\ldots,n\}$, set $r_j\defeq1$ if $\lambda_j=0$ and $r_j\defeq\abs{\lambda_j}$ otherwise. Then $(\forall\,j\in\{1,\ldots,n\})[r_j>0\land r_jd_j=\lambda_j]$.\\
Set $S\defeq U\operatorname{diag}(\sqrt{r_1},\ldots,\sqrt{r_n})$. Then $S$ is invertible and $SDS^*=U\operatorname{diag}(r_1d_1,\ldots,r_nd_n)U^*=A$.
\end{enumerate}
\hlt{Structural Commentary}: Spectral diagonalization separates phase from positive radius, which the congruence factor absorbs.\\
\hlt{Key Insight}: Radial square roots provide the diagonal scaling in $S$.
\end{proof}
\end{proposition}

\begin{perspective} \label{per:hermitian-inertia-as-real-phase-inertia} \hlt{Hermitian Inertia as Real Phase Inertia}\\
Proposition \ref{prop:normal-matrix-phase-inertia-form} reduces a normal matrix to its eigenvalue phases under $*$-congruence.\\
Hermitian eigenvalues are real by Theorem \ref{thm:self-adjoint-eigenvalues-real}, so their nonzero phases are $\pm1$;\\
Theorem \ref{thm:sylvester-inertia-theorem-hermitian} strengthens this reduction to a complete classification by sign multiplicities.
\end{perspective}
